% GENERATED FILE. Edit canonical lessons, then run npm run generate:books.
% canonical-source-hash: fnv1a64:7d945ffc7a764565
% canonical-lessons: MR-C162-petavne, MR-C162-vijhavne, MR-C162-udi-marne, MR-C162-bolavne, MR-C162-abhar-manne

\chapter{To light, To put out, To jump, To call over, To thank}
\label{ch:mr-to-light-to-put-out-to-jump-to-call-over-to-thank}
\hlchaptermodality{162}

\begin{chapteropening}
\textbf{By the end of this chapter:} \emph{I can say to light, to put out, to jump, to call over and to thank.}

Complete the last lesson of chapter 162: I can say to light, to put out, to jump, to call over and to thank.
\end{chapteropening}

\section[\texorpdfstring{\mr{पेटवणे} (petavṇe)}{पेटवणे (petavṇe)}]{\mr{पेटवणे} (petavṇe) — to light, to set alight}

\label{lesson:MR-C162-petavne}



\begin{quote}
Before the new one: say the Marathi for an advert, then the Marathi for a receipt.
\end{quote}

\subsection*{You'll want to know: \mr{पेटवणे}}
\textbf{\mr{पेटवणे}} — \emph{petavṇe} — \textquotedblleft{}to light, to set alight\textquotedblright{}.

\begin{grammarlens}[title={What you've built}]
One more word for this chapter.
\end{grammarlens}

\subsection*{Your turn}
\begin{itemize}
\raggedright
  \item \emph{Say it:} \emph{petavṇe}
  \item \emph{Say it:} \emph{petavṇe}, once more
  \item \emph{Say it:} say \emph{petavṇe}
  \item \emph{Recall:} say the Marathi for an advert, then the Marathi for a receipt, then say \emph{petavṇe} again
\end{itemize}

\begin{tcolorbox}[breakable,colback=teal!4,colframe=teal!35!black,title={Before you move on}]
What does \mr{पेटवणे} mean? (\textquotedblleft{}To light, to set alight\textquotedblright{}.) Say it once more.
\end{tcolorbox}
\section[\texorpdfstring{\mr{विझवणे} (vijhavṇe)}{विझवणे (vijhavṇe)}]{\mr{विझवणे} (vijhavṇe) — to put out}

\label{lesson:MR-C162-vijhavne}



\begin{quote}
Before the new one: say the Marathi for a receipt, then the Marathi for to light.
\end{quote}

\subsection*{You'll want to know: \mr{विझवणे}}
\textbf{\mr{विझवणे}} — \emph{vijhavṇe} — \textquotedblleft{}to put out\textquotedblright{}.

\begin{grammarlens}[title={What you've built}]
One more word for this chapter.
\end{grammarlens}

\subsection*{Your turn}
\begin{itemize}
\raggedright
  \item \emph{Say it:} \emph{vijhavṇe}
  \item \emph{Say it:} \emph{vijhavṇe}, once more
  \item \emph{Say it:} say \emph{vijhavṇe}
  \item \emph{Recall:} say the Marathi for a receipt, then the Marathi for to light, then say \emph{vijhavṇe} again
\end{itemize}

\begin{tcolorbox}[breakable,colback=teal!4,colframe=teal!35!black,title={Before you move on}]
What does \mr{विझवणे} mean? (\textquotedblleft{}To put out\textquotedblright{}.) Say it once more.
\end{tcolorbox}
\section[\texorpdfstring{\mr{उडी} \mr{मारणे} (uḍī mārṇe)}{उडी मारणे (uḍī mārṇe)}]{\mr{उडी} \mr{मारणे} (uḍī mārṇe) — to jump}

\label{lesson:MR-C162-udi-marne}



\begin{quote}
Before the new one: say the Marathi for to light, then the Marathi for to put out.
\end{quote}

\subsection*{You'll want to know: \mr{उडी} \mr{मारणे}}
\textbf{\mr{उडी} \mr{मारणे}} — \emph{uḍī mārṇe} — \textquotedblleft{}to jump\textquotedblright{}.

\begin{grammarlens}[title={What you've built}]
One more word for this chapter.
\end{grammarlens}

\subsection*{Your turn}
\begin{itemize}
\raggedright
  \item \emph{Say it:} \emph{uḍī mārṇe}
  \item \emph{Say it:} \emph{uḍī mārṇe}, once more
  \item \emph{Say it:} say \emph{uḍī mārṇe}
  \item \emph{Recall:} say the Marathi for to light, then the Marathi for to put out, then say \emph{uḍī mārṇe} again
\end{itemize}

\begin{tcolorbox}[breakable,colback=teal!4,colframe=teal!35!black,title={Before you move on}]
What does \mr{उडी} \mr{मारणे} mean? (\textquotedblleft{}To jump\textquotedblright{}.) Say it once more.
\end{tcolorbox}
\section[\texorpdfstring{\mr{बोलावणे} (bolāvṇe)}{बोलावणे (bolāvṇe)}]{\mr{बोलावणे} (bolāvṇe) — to call over, to invite}

\label{lesson:MR-C162-bolavne}



\begin{quote}
Before the new one: say the Marathi for to put out, then the Marathi for to jump.
\end{quote}

\subsection*{You'll want to know: \mr{बोलावणे}}
\textbf{\mr{बोलावणे}} — \emph{bolāvṇe} — \textquotedblleft{}to call over, to invite\textquotedblright{}.

\begin{grammarlens}[title={What you've built}]
One more word for this chapter.
\end{grammarlens}

\subsection*{Your turn}
\begin{itemize}
\raggedright
  \item \emph{Say it:} \emph{bolāvṇe}
  \item \emph{Say it:} \emph{bolāvṇe}, once more
  \item \emph{Say it:} say \emph{bolāvṇe}
  \item \emph{Recall:} say the Marathi for to put out, then the Marathi for to jump, then say \emph{bolāvṇe} again
\end{itemize}

\begin{tcolorbox}[breakable,colback=teal!4,colframe=teal!35!black,title={Before you move on}]
What does \mr{बोलावणे} mean? (\textquotedblleft{}To call over, to invite\textquotedblright{}.) Say it once more.
\end{tcolorbox}
\section[\texorpdfstring{\mr{आभार} \mr{मानणे} (ābhār mānṇe)}{आभार मानणे (ābhār mānṇe)}]{\mr{आभार} \mr{मानणे} (ābhār mānṇe) — to thank}

\label{lesson:MR-C162-abhar-manne}



\begin{quote}
Before the new one: say the Marathi for to jump, then the Marathi for to call over.
\end{quote}

\subsection*{You'll want to know: \mr{आभार} \mr{मानणे}}
\textbf{\mr{आभार} \mr{मानणे}} — \emph{ābhār mānṇe} — \textquotedblleft{}to thank\textquotedblright{}.

\begin{grammarlens}[title={What you've built}]
One more word for this chapter.
\end{grammarlens}

\subsection*{Your turn}
\begin{itemize}
\raggedright
  \item \emph{Say it:} \emph{ābhār mānṇe}
  \item \emph{Say it:} \emph{ābhār mānṇe}, once more
  \item \emph{Say it:} say \emph{ābhār mānṇe}
  \item \emph{Recall:} say the Marathi for to jump, then the Marathi for to call over, then say \emph{ābhār mānṇe} again
\end{itemize}

\begin{tcolorbox}[breakable,colback=teal!4,colframe=teal!35!black,title={Before you move on}]
What does \mr{आभार} \mr{मानणे} mean? (\textquotedblleft{}To thank\textquotedblright{}.) Say it once more.
\end{tcolorbox}
